\documentclass[../../main.tex]{subfiles}
\begin{document}

\begin{flushright}
\footnotesize\textit{【自動文字起こし・要確認】}
\end{flushright}

\noindent
{\bf [解]} 点 $X$ に対し，$\vec{AX} = \vec{x}$ とおく．

\bigskip

\paragraph{(1)}

$AB, AC$ 上に，端点除いて奇数個の格子点がある時，$(a,b), (c,d)$ を互いに素な自然数として
\begin{align}
\vec{b} &= 2k \begin{pmatrix} a \\ b \end{pmatrix}, \\
\vec{c} &= 2l \begin{pmatrix} c \\ d \end{pmatrix} \quad (k, l \in \mathbb{N})
\end{align}
とおいて良い．(点 $A, B, C$ に関する対称性)
\begin{align}
\vec{BC} &= 2l \begin{pmatrix} c \\ d \end{pmatrix} - 2k \begin{pmatrix} a \\ b \end{pmatrix} \\
&= 2 \begin{pmatrix} lc - ka \\ ld - kb \end{pmatrix}
\end{align}
において，$lc - ka, ld - kb \in \mathbb{Z}$ より，たしかに $BC$ 上には，端点を除いて奇数個の格子点がある．

\bigskip

\paragraph{(2)}

(1)と同じく，
\begin{align}
\vec{b} &= 4 \begin{pmatrix} a \\ b \end{pmatrix}, \\
\vec{c} &= 4 \begin{pmatrix} c \\ d \end{pmatrix}
\end{align}
とおけるから，$\triangle ABC$ の面積 $S$ として
\begin{align}
S = \dfrac{1}{2} | 16 ad - 16 bc | = 8 |ad - bc|
\end{align}
$ad - bc \in \mathbb{Z}$ より，たしかに $S$ は 8 の倍数．

\begin{figure}[htb]
\centering
\begin{tikzpicture}[scale=0.8, >=stealth]
\coordinate (A) at (0,0);
\coordinate (B) at (2,3);
\coordinate (C) at (4,1);
\draw[thick] (A) -- (B) -- (C) -- cycle;
\filldraw (A) circle (2pt) node[below] {$A$};
\filldraw (B) circle (2pt) node[above] {$B$};
\filldraw (C) circle (2pt) node[right] {$C$};
\foreach \i in {1,2,3} {
  \filldraw ($(A)!\i/4!(B)$) circle (1.5pt);
  \filldraw ($(A)!\i/4!(C)$) circle (1.5pt);
}
\end{tikzpicture}
\caption{$\triangle ABC$ の辺 $AB, AC$ 上の格子点（各辺を4等分する点）}
\end{figure}

\newpage

\end{document}
